% Figure: the SW07 policy counterfactual sized to offset the GDP trough.
% Data: figs/data/us_cf.csv (app/scripts/paper_figs/us_spread_cf.mjs).
\begin{figure}[t]
\centering
\begin{tikzpicture}
\begin{groupplot}[
  group style={group size=2 by 1, horizontal sep=1.6cm},
  paperpanel, width=0.46\textwidth, height=0.32\textwidth,
]
\nextgroupplot[title={Real GDP (\% deviation from baseline)},
  legend style={font=\scriptsize, draw=none, fill=none, at={(0.95,0.05)}, anchor=south east}]
\addplot[zeroline] coordinates {(1,0)(20,0)};
\addplot[trueline] table[x=q, y=gdp_scen_lvl, col sep=comma] {figs/data/us_cf.csv};
\addplot[cfline]   table[x=q, y=gdp_cf_lvl,   col sep=comma] {figs/data/us_cf.csv};
\legend{spread scenario, with policy response}
\nextgroupplot[title={Policy rate (pp deviation from baseline)}]
\addplot[zeroline] coordinates {(1,0)(20,0)};
\addplot[trueline] table[x=q, y=pol_scen, col sep=comma] {figs/data/us_cf.csv};
\addplot[cfline]   table[x=q, y=pol_cf,   col sep=comma] {figs/data/us_cf.csv};
\end{groupplot}
\end{tikzpicture}
\caption{What should policy do? The credit-spread scenario (solid) and the same scenario with
an unanticipated monetary-policy shock from the Smets--Wouters model layered on top (dashed),
sized so that the easing offsets the output trough. Left: the level of real GDP relative to
baseline; the counterfactual closes the gap the shock opened. Right: the policy-rate paths. The
scenario's own rate falls because the model expects policy to respond; the counterfactual adds
the deliberate surprise easing on impact, which then decays along the model's policy rule.}
\label{fig:cf}
\end{figure}
